\section{Results}

\subsection{Forecast Quality on Overlapping Test Forecasts}

Table~\ref{tab:primary-comparison} reports the three canonical models central to the paper. The primary hybrid has the highest all-overlap Empty AUPRC point estimate (0.8514), followed closely by the historical average (0.8497); LightGBM reaches 0.8382. The hybrid also has the lowest Brier score among these three models (0.0928). These forecast results cover 388,032 overlapping rows from 4,042 anchors. Appendix~A retains the extended seven-model comparison and explains the supplementary status of the test-ranked balanced hybrid.

\begin{table*}[!t]
\caption{Canonical three-model forecast comparison on all overlapping held-out test forecasts (Empty $=1$). Scores are uncalibrated. Precision, recall, and F1 use a descriptive 0.5 score cutoff; Table~\ref{tab:policy-comparison} uses separately selected validation policy thresholds.}
\label{tab:primary-comparison}
\centering
\small
\begin{tabular}{lccccc}
\toprule
Model & Empty AUPRC & Empty Brier & Empty precision & Empty recall & Empty F1 \\
\midrule
Historical average & 0.8497 & 0.0974 & 0.7264 & 0.8015 & 0.7621 \\
LightGBM & 0.8382 & 0.0956 & 0.7610 & 0.7593 & 0.7601 \\
Primary hybrid & 0.8514 & 0.0928 & 0.7677 & 0.7556 & 0.7616 \\
\bottomrule
\end{tabular}
\end{table*}


The strong schedule baseline is a substantive result: weekday and time-of-day structure carries much of the Empty-class predictive signal in this building. The hybrid's point AUPRC margins are 0.0017 over the historical baseline and 0.0132 over LightGBM. Section~V-C evaluates the corresponding daily-horizon differences.

\subsection{Validation-Selection Stability}

The canonical weight vector attains the nominal grid maximum of 0.728638 validation AUPRC. The runner-up is 0.000032 lower, and eight of the 231 legal vectors lie within 0.0005. Table~\ref{tab:selection-stability} reports separate validation-only perturbation diagnostics on 39 non-overlapping validation policy horizons. Nearby blends change rank under validation-horizon perturbation, and threshold selection varies across daily-block bootstrap resamples. The canonical vector is therefore one validation-selected solution within a competitive neighborhood.

\begin{table*}[!t]
\caption{Validation-only selection-stability diagnostics. The blend perturbation uses 39 non-overlapping validation policy horizons and the top 12 full-overlap grid candidates; the canonical all-overlap selection remains primary. Threshold diagnostics apply the declared validation policy rule.}
\label{tab:selection-stability}
\centering
\scriptsize
\setlength{\tabcolsep}{2pt}
\begin{tabular}{p{0.26\textwidth}p{0.40\textwidth}r}
\toprule
Diagnostic & Result & Frequency / value \\
\midrule
Full-overlap 231-weight grid & Nominal rank-1 AUPRC; rank-1 minus rank-2; candidates within 0.0005 & 0.728638; 0.000032; 8 \\
Daily-block bootstrap, top-12 blends (1,000) & Canonical $(0.15,0.60,0.25)$ selected; most frequent $(0.05,0.80,0.15)$ & 23 (2.3\%); 365 (36.5\%) \\
Leave-one-validation-horizon-out, top-12 blends (39) & Canonical selected; most frequent $(0.05,0.75,0.20)$ & 1 (2.6\%); 31 (79.5\%) \\
Daily-block bootstrap, primary threshold (1,000) & Canonical $\tau=0.875$; $\tau=0.850$; $\tau=0.950$ & 111 (11.1\%); 376 (37.6\%); 375 (37.5\%) \\
Leave-one-validation-horizon-out, primary threshold (39) & Canonical $\tau=0.875$ & 36 (92.3\%) \\
\bottomrule
\end{tabular}
\end{table*}


\subsection{Validation-Selected Conflict and Load-Opportunity Outcomes}

Table~\ref{tab:policy-comparison} and Fig.~\ref{fig:policy} report the policy scope: 43 non-overlapping test horizons and 4,128 15-min target intervals. The horizons are anchored by midnight-labelled completed input bins and have an effective 00:15 availability boundary. Every listed threshold met the empirical validation $10\%$ interval-conflict condition. The primary hybrid used $\tau=0.875$ after a validation conflict rate of 8.75\%. On test, it recommended 259 intervals. All 259 were subsequently camera-label-empty. Those intervals formed 14 recommended windows, all camera-label-safe under the observed label, and coincided with 490.1~kWh of \safeopportunity{} (11.40~kWh/horizon-day).

\begin{table*}[!t]
\caption{Validation-selected one-hour stable-window policies evaluated on 43 non-overlapping test horizons. Anchors are midnight-labelled completed input bins with effective 00:15 boundaries. Every threshold maximizes validation safe opportunity subject to empirical validation interval conflict $\leq10\%$. ``Safe'' means recommended and subsequently camera-label-empty; opportunity is an offline processed-load-proxy overlap.}
\label{tab:policy-comparison}
\centering
\footnotesize
\setlength{\tabcolsep}{3pt}
\begin{tabular}{lccccccc}
\toprule
Model & $\tau$ & Val. conflict & Test conflict & Safe opportunity & Safe/rec. intervals & Safe/rec. windows & kWh/day \\
\midrule
Historical average & 0.950 & 8.33\% & 0.00\% (0/66) & 94.6 & 66/66 & 12/12 & 2.20 \\
LightGBM & 0.950 & 9.61\% & 4.15\% (11/265) & 493.9 & 254/265 & 16/19 & 11.49 \\
Primary hybrid & 0.875 & 8.75\% & 0.00\% (0/259) & 490.1 & 259/259 & 14/14 & 11.40 \\
\bottomrule
\end{tabular}
\end{table*}


LightGBM produced a nearly equal 493.9~kWh opportunity, with 11 of 265 recommended intervals camera-labelled occupied (4.15\%) and three of 19 recommended windows containing an occupancy conflict. The historical average was also conflict-free under this label and coincided with 94.6~kWh of processed load-proxy opportunity. Figure~\ref{fig:policy} compares these three pre-specified models.

Taken together, the forecast and policy results explain why the hybrid is the primary case-study model. It preserves the strong schedule baseline's Empty-class discrimination while producing the highest nominal AUPRC of the three models. At its validation-selected threshold, it identifies nearly the same held-out load-opportunity proxy as LightGBM (490.1 versus 493.9~kWh) with fewer observed camera-label conflicts (0 versus 11), and substantially more opportunity than the conflict-free historical policy (490.1 versus 94.6~kWh). These results give the hybrid the strongest observed opportunity--conflict trade-off in the saved test period.

\begin{figure*}[!t]
  \centering
  \includegraphics[width=\textwidth]{fig_policy_comparison.png}
  \caption{Fixed held-out policy outcomes for the schedule baseline, LightGBM, and the validation-selected primary hybrid. Thresholds were selected on 39 validation daily horizons to maximize offline safe load opportunity subject to an empirical interval conflict rate of at most 10\%; the plotted values come from 43 post-bin test horizons. ``Safe'' denotes recommended and subsequently camera-label-empty.}
  \label{fig:policy}
\end{figure*}

\subsection{Conditional Uncertainty and Diagnostics}

The daily-block uncertainty results in Table~\ref{tab:uncertainty} condition on the previously selected weights and thresholds. The primary hybrid's daily-horizon AUPRC is 0.8522 with a 95\% percentile interval of [0.7918, 0.9016]. Its aggregation unit differs from the all-overlap value in Table~\ref{tab:primary-comparison}. The paired AUPRC intervals for primary minus LightGBM and primary minus historical average both cross zero. The paired safe-opportunity interval crosses zero for primary minus LightGBM; the primary-minus-historical interval is positive. Figure~\ref{fig:uncertainty} displays these paired contrasts on zero-centered axes.

\begin{table*}[!t]
\caption{Conditional daily-block uncertainty from 2,000 paired bootstrap resamples of 43 non-overlapping test horizons. Daily-horizon AUPRC uses a different aggregation from the all-overlap AUPRC in Table~\ref{tab:primary-comparison}. All intervals condition on selected weights and thresholds.}
\label{tab:uncertainty}
\centering
\small
\begin{tabular}{llrrr}
\toprule
Estimate type & Quantity & Point estimate & 95\% low & 95\% high \\
\midrule
Primary model & Daily-horizon Empty AUPRC & 0.8522 & 0.7918 & 0.9016 \\
Primary model & Safe opportunity (kWh) & 490.1 & 153.1 & 966.7 \\
Paired contrast & Primary $-$ LightGBM AUPRC & +0.0163 & -0.0050 & +0.0426 \\
Paired contrast & Primary $-$ historical AUPRC & +0.0042 & -0.0304 & +0.0377 \\
Paired contrast & Primary $-$ LightGBM opportunity (kWh) & -3.8 & -208.2 & +174.6 \\
Paired contrast & Primary $-$ historical opportunity (kWh) & +395.5 & +127.0 & +786.2 \\
\bottomrule
\end{tabular}
\end{table*}


\begin{figure*}[!t]
  \centering
  \includegraphics[width=\textwidth]{fig_paired_uncertainty.png}
  \caption{Paired 95\% daily-block bootstrap intervals (2,000 resamples) for the primary hybrid minus each canonical baseline. The dashed line marks zero. Intervals condition on fixed, validation-selected model weights and operating thresholds.}
  \label{fig:uncertainty}
\end{figure*}

The primary hybrid's held-out score diagnostics are Brier score 0.0928, log loss 0.2984, and 10-bin quantile ECE 0.0253. At the fixed primary threshold, a one-hour or two-hour minimum window retains 490.1~kWh and zero observed camera-label conflicts; a four-hour requirement retains 407.8~kWh and 11 fully camera-label-safe windows. These sensitivity results preserve the selected threshold. Every observed primary-hybrid camera-label conflict block is zero, so a nonparametric bootstrap also returns a zero conflict interval. Under an independent-trial approximation, zero events yield one-sided 95\% upper bounds of 1.15\% for 259 intervals and 19.3\% for 14 windows. The clustering of intervals and windows makes these bounds descriptive.
